
\subsubsection{6.1 Summary of Key Findings}

\textbf{The domain content differentiation result is the primary empirical contribution.} $\eta^2$ > 0.91 means that domain membership explains over 91\% of total variance in each cognitive content proxy across 4,200 documents. Prior data mixing research (RegMix, DoReMi, DCLM) has assumed this differentiation without measuring it directly. The h-m1 result provides the first quantitative evidence. The caveat that generated texts were used (not real Pile documents) is acknowledged; directional claims are robust.

\textbf{The trajectory non-uniformity result establishes the prerequisite for within-family panel analysis.} The infrastructure validation --- \texttt{doc\textbackslash \_idx.npy} identity mapping confirmed for 134M entries, JSONL.zst streaming completing in 3.5 minutes for 600,000 documents, 22 $\times$ 154 trajectory matrix produced without NaN --- constitutes the missing methodological link between Pythia's training procedure and domain exposure analysis. Cross-scale Spearman $\rho$ = 1.0 confirms that trajectories are a property of data ordering.

\textbf{The 70M reversal warrants replication before interpretation.} The competing explanations, ranked by plausibility, are: (1) MMLU floor effect at 70M --- MMLU is near-random ($\approx 25$\%) at 70M scale, making $\rho$(anything, MMLU) noise-dominated [HIGH]; (2) N = 10 checkpoint sampling artifact with non-uniform step distribution [HIGH]; (3) Pile-CC multicollinearity confounding Wikipedia's univariate Spearman [MEDIUM]; (4) genuine domain-benchmark reversal at all scales [LOW]. The 1B and 6.9B evaluations with full 154-checkpoint coverage will adjudicate. The 70M result is reported faithfully but cannot be treated as informative about the underlying domain-benchmark relationship.

\textbf{The Books3 zero-exposure finding is a concrete, reproducible infrastructure result.} Six Pile domains (Books3, OpenWebText2, GitHub, OpenSubtitles, BookCorpus2, YoutubeSubtitles) have zero exposure in the first-shard sample. This is consistent with The Pile's block structure concentrating these domains in later shards. The same pattern affecting six distinct domains across all three model sizes confirms this is a data structure property, not a pipeline error. This finding constitutes a reproducibility warning for future work using partial Pile samples.

\subsubsection{6.2 Limitations}

\textbf{L1: Six Pile domains have zero exposure in the first-shard PoC (Books3, OpenWebText2, GitHub, OpenSubtitles, BookCorpus2, YoutubeSubtitles).} This blocks both P2 (Books3 $\rightarrow$ HellaSwag specificity) and the formal-syntax component of the causal mechanism (GitHub). Fix: run \texttt{build\textbackslash \_lookup\textbackslash \_direct.py} across all 30 Pile shards (estimated ~150 compute-hours, ~330MB output). The 10 high-variance domains including Wikipedia remain testable with the current lookup.

\textbf{L2: Benchmark evaluation cache was incomplete at analysis time (70M: 10/154; 1B: 2/154; 6.9B: 0/154).} This prevents reliable Spearman analysis (pre-registered minimum: N = 100) and panel regression (minimum 3 entities). The evaluation framework was confirmed correct and running on $5\times$ H100 NVL GPUs at reporting time. Full completion enables re-running all h-m2 and h-m3 analyses without methodological changes.

\textbf{L3: h-m1 used domain-representative generated texts rather than actual Pile documents.} Effect sizes ($\eta^2$ > 0.91) may be inflated. Generated texts likely have cleaner domain-typicality than actual mixed-content Pile documents. A background validation run on real Pile documents (1,000 documents $\times$ 22 domains via \texttt{monology/pile-uncopyrighted}) was running in parallel but did not complete within the execution window. Directional claims (Wikipedia highest entity density; BookCorpus2 highest narrative coherence; GitHub highest formal syntax density) are robust.

\textbf{L4: Observational design --- correlations do not establish causal domain effects on benchmark performance.} Within-family variation controls for architecture and data composition confounds, but is not equivalent to a controlled intervention. The panel fixed effects control for scale, but unobserved confounders at the checkpoint level cannot be ruled out.

\subsubsection{6.3 Connection to Prior Work}

The domain content differentiation result ($\eta^2$ > 0.91) empirically grounds the assumption underlying Data Mixing Laws \citep{Ye2024} and RegMix \citep{Liu2024}: that domains are qualitatively distinct in the features relevant to downstream performance. Both prior works treat this as a given; we measure it.

The trajectory extraction infrastructure extends Pythia \citep{Biderman2023}: the Pythia paper documents the architecture and checkpoints; we exploit the block data structure to derive time-varying domain exposure covariates, an application not previously demonstrated.

The preliminary 70M reversal is partially consistent with CoLoR-Filter \citep{Brandfonbrener2024}: tasks require different data, but Wikipedia's role at small scale may be more general-purpose than MMLU-specific.

The Books3 zero-exposure finding is consistent with BiMix \citep{Ge2024}: minimum domain exposure thresholds must be exceeded before domain signals become detectable.

\subsubsection{6.4 Future Work}

The most critical next step is building the full 134M-document domain lookup across all 30 Pile shards to obtain non-zero Books3 and GitHub exposure signals. With the complete evaluation cache (462 evaluations) and full lookup, the h-m2 Spearman analysis and h-m3 panel regression can be executed without modification.

Two alternative hypotheses raised by the 70M reversal merit investigation: (1) whether domain-benchmark specificity is scale-emergent, appearing only at $\geq 1B$ parameters where MMLU performance clears floor levels; and (2) whether early-phase training (steps 0--10K) shows general improvement from any coherent text, masking later-phase domain-specific alignment --- testable by splitting the 154 checkpoints into early and late phases.

Additionally, replacing Books3 with a high-variance focal domain (e.g., StackExchange, std = 0.01496) in the specificity test would enable P2-equivalent analysis with the current data.

